% (c) Nikita Lisitsa, lisyarus@gmail.com, 2024

\documentclass[10pt]{beamer}

\usepackage[T2A]{fontenc}
\usepackage[russian]{babel}
\usepackage{minted}

\usepackage{graphicx}
\graphicspath{ {./images/} }

\usepackage{adjustbox}

\usepackage{color}
\usepackage{soul}

\usepackage{hyperref}

\usetheme{metropolis}

\definecolor{red}{rgb}{1,0,0}
\definecolor{green}{rgb}{0,0.5,0}
\definecolor{blue}{rgb}{0,0,1}
\definecolor{codebg}{RGB}{29,35,49}
\definecolor{lightbg}{RGB}{253,246,227}
\setminted{fontsize=\footnotesize}

\makeatletter
\newcommand{\slideimage}[1]{
  \begin{figure}
    \begin{adjustbox}{width=\textwidth, totalheight=\textheight-2\baselineskip-2\baselineskip,keepaspectratio}
      \includegraphics{#1}
    \end{adjustbox}
  \end{figure}
}
\makeatother

\title{Фотореалистичный рендеринг\quad\quad\quad\quad\quad\quad \textit{(aka raytracing)}}
\subtitle{Лекция 6: Alias method, bidirectional path tracing, Metropolis light transport, формат glTF}
\date{2024}

\setbeamertemplate{footline}[frame number]

\begin{document}

\frame{\titlepage}

\begin{frame}
\frametitle{Importance sampling lights}
\begin{itemize}
\item При importance sampling мы выбирали случайный источник равновероятно
\pause
\item Это не всегда самая лучшая стратегия:
\pause
\begin{itemize}
\item Источники отличаются размером
\pause
\item Источники отличаются яркостью
\end{itemize}
\pause
\item \begin{math}\Longrightarrow\end{math} Хочется какую-то эвристику
\end{itemize}
\end{frame}

\begin{frame}
\frametitle{Weighted lights sampling}
\begin{itemize}
\item В качестве эвристического веса \begin{math}w_i\end{math} источника света можно взять
\pause
\begin{itemize}
\item Яркость источника
\pause
\item Площадь поверхности источника
\pause
\item Их произведение
\end{itemize}
\pause
\item Вероятность выбрать конкретный источник равна
\begin{equation*}p_i = \frac{w_i}{\sum w_j}\end{equation*}
\end{itemize}
\end{frame}

\begin{frame}
\frametitle{Uniform lights sampling}
\slideimage{uniform_lights.png}
\end{frame}

\begin{frame}
\frametitle{Weighted lights sampling}
\slideimage{weighted_lights.png}
\end{frame}

\begin{frame}
\frametitle{Weighted lights sampling}
\begin{itemize}
\item Для семплинга источников света с учётом такой эвристики, нужно уметь семплить произвольное дискретное распределение, заданное набором вероятностей
\pause
\item Эта операция выполняется почти на каждый луч \begin{math}\Longrightarrow\end{math} хочется делать это быстро
\end{itemize}
\end{frame}

\begin{frame}
\frametitle{Weighted sampling: постановка задачи}
\begin{itemize}
\item Есть дискретное распределение с заданными вероятностями исходов \begin{math}p_i\end{math} для \begin{math}i \in \{1\dots N\}\end{math}
\pause
\item Хотим научиться семплить случайную величину \begin{math}X\end{math} такую, что \begin{math}\mathbb{P}(X = i) = p_i\end{math}
\end{itemize}
\end{frame}

\begin{frame}
\frametitle{Weighted sampling: линейный алгоритм}
\begin{itemize}
\item Генерируем случайную величину \begin{math}U \sim U(0, 1)\end{math}
\pause
\item Если \begin{math}U \leq p_1\end{math}, возвращаем \begin{math}X = 1\end{math}
\pause
\item Если \begin{math}p_1 < U \leq p_1 + p_2\end{math}, возвращаем \begin{math}X = 2\end{math}
\pause
\item Если \begin{math}p_1 + p_2 < U \leq p_1 + p_2 + p_3\end{math}, возвращаем \begin{math}X = 3\end{math}
\pause
\item И т.д.
\end{itemize}
\end{frame}

\begin{frame}
\frametitle{Weighted sampling: линейный алгоритм}
\begin{itemize}
\item Есть разные варианты реализации:
\pause
\begin{itemize}
\item Можно считать сумму вероятностей всех посещённых индексов в процессе семплинга
\pause
\item Можно заранее предподсчитать префиксные суммы \begin{math}s_k = \sum\limits_{i=1}^k p_i\end{math}
\pause
\item Можно вычитать вероятности посещённых индексов из \begin{math}U\end{math}
\end{itemize}
\end{itemize}
\end{frame}

\begin{frame}
\frametitle{Weighted sampling: логарифмический алгоритм}
\usemintedstyle{solarized-light}
\begin{itemize}
\item Линейная сложность -- очень медленно, если источников света много
\pause
\item Если заранее предподсчитать префиксные суммы, то нужный индекс можно найти бинарным поиском
\pause
\item Находим минимальный \begin{math}i\end{math} такой, что \begin{math}U \leq s_i\end{math}
\pause
\item В C++ это \mintinline{cpp}|std::upper_bound|
\pause
\item Сложность такого подхода -- логарифм от числа источников света
\end{itemize}
\end{frame}

\begin{frame}
\frametitle{Weighted sampling: alias method}
\usemintedstyle{solarized-light}
\begin{itemize}
\item Логарифм -- уже очень хорошо, но не идеально
\pause
\begin{itemize}
\item Сложность всё ещё растёт с числом источников света
\pause
\item Бинарный поиск плохо дружит с кешами памяти, так как делает много практически случайных обращений к массиву
\end{itemize}
\pause
\item Есть метод, работающий за константное время с линейной предобработкой: \textit{alias method}
\end{itemize}
\end{frame}

\begin{frame}
\frametitle{Alias method (Walker, 1974)}
\usemintedstyle{solarized-light}
\begin{itemize}
\item Рассмотрим особый класс дискретных распределений, заданных двумя массивами \begin{math}q_i \in [0, 1]\end{math} \textit{(probability table)} и \begin{math}k_i \in \{1\dots N\}\end{math} \textit{(alias table)} и алгоритмом семплинга:
\pause
\begin{itemize}
\item Генерируем равномерно \begin{math}i \in \{1\dots N\}\end{math} и \begin{math}u \in [0, 1]\end{math}
\pause
\item Если \begin{math}u \leq q_i\end{math}, возвращаем \begin{math}i\end{math}
\pause
\item Иначе, возвращаем \begin{math}k_i\end{math}
\end{itemize}
\pause
\item Утверждение: \textbf{любое} дискретное распределение можно привести к такому виду
\pause
\item Элемент \begin{math}k_i\end{math} называется \textit{alias'ом} элемента \begin{math}i\end{math}
\end{itemize}
\end{frame}

\begin{frame}
\frametitle{Alias method: доказательство}
\usemintedstyle{solarized-light}
\begin{itemize}
\item 1. Положим \begin{math}q_i = N \cdot p_i\end{math} и \begin{math}k_i = i\end{math}
\pause
\item Заметим, что \begin{math}\sum q_i = N \cdot \sum p_i = N\end{math}
\pause
\item 2. Если все \begin{math}q_i = 1\end{math}, то все \begin{math}p_i = \frac{1}{N}\end{math}, и алгоритм можно завершить
\pause
\item 3. Иначе, выберем \begin{math}q_i < 1\end{math} и \begin{math}q_j > 1\end{math}
\pause
\begin{itemize}
\item Они существуют из-за того, что \begin{math}\sum q_i = N\end{math}
\end{itemize}
\pause
\item 4. `Откусим' часть вероятности \begin{math}q_j\end{math}, чтобы дополнить ей \begin{math}q_i\end{math}:
\pause
\begin{itemize}
\item 4.1. Положим \begin{math}k_i = j\end{math}
\pause
\item 4.2. Обновим \begin{math}q_j \leftarrow q_j - (1 - q_i)\end{math}
\end{itemize}
\pause
\item 5. Удалим \begin{math}i\end{math}-ый элемент из рассмотрения
\pause
\item Сумма \begin{math}q_i\end{math} для оставшихся элементов равна \begin{math}N - q_i - (1 - q_i) = N - 1\end{math}
\pause
\item 6. \begin{math}\Longrightarrow\end{math} Инвариант на сумму \begin{math}q_i\end{math} сохранился, возвращаемся к пункту 2
\end{itemize}
\end{frame}

\begin{frame}
\frametitle{Alias method: алгоритм}
\usemintedstyle{solarized-light}
\begin{itemize}
\item Приведённое доказательство конструктивно, и уже является алгоритмом построения массивов \begin{math}q_i\end{math} и \begin{math}k_i\end{math}
\pause
\item Однако, на каждом шаге нам нужно за линейное время искать элементы \begin{math}q_i < 1\end{math} и \begin{math}q_j > 1\end{math} \begin{math}\Longrightarrow\end{math} квадратичная сложность предподсчёта
\pause
\item К тому же, алгоритм страдает от ошибок округления: инвариант \begin{math}\sum q_i = N\end{math} может нарушиться, и один из элементов \begin{math}q_i, q_j\end{math} может не найтись
\pause
\item Вместо линейного поиска можно хранить две очереди с элементами меньше и больше 1
\pause
\item Можно заменить вычисление \begin{math}q_j - (1 - q_i)\end{math} численно более стабильным \begin{math}(q_j + q_i) - 1\end{math}
\pause
\item \begin{math}\Longrightarrow\end{math} \textit{Vose alias method}
\end{itemize}
\end{frame}

\begin{frame}
\frametitle{Vose alias method (Vose, 1991)}
\usemintedstyle{solarized-light}
\begin{itemize}
\item 1. Положим \begin{math}q_i = N \cdot p_i\end{math}
\pause
\item 2. Заведём две очереди:
\pause
\begin{itemize}
\item 2.1. \begin{math}\operatorname{small}\end{math} с индексами, для которых \begin{math}q_i < 1\end{math}
\pause
\item 2.2. \begin{math}\operatorname{large}\end{math} с индексами, для которых \begin{math}q_i \geq 1\end{math}
\end{itemize}
\pause
\item 3. Пока обе очереди не пусты:
\pause
\begin{itemize}
\item 3.1. Достанем элементы \begin{math}i \in \operatorname{small}\end{math} и \begin{math}j \in \operatorname{large}\end{math}
\pause
\item 3.2. Положим \begin{math}k_i = j\end{math}
\pause
\item 3.3. Обновим \begin{math}q_j \leftarrow (q_i + q_j) - 1\end{math}
\pause
\item 3.4. Если \begin{math}q_j < 1\end{math}, положим \begin{math}j\end{math} в \begin{math}\operatorname{small}\end{math}, иначе -- в \begin{math}\operatorname{large}\end{math}
\end{itemize}
\pause
\item 4. Если одна из очередей опустела, то для всех элементов второй очереди положим \begin{math}q_i = 1\end{math} и \begin{math}k_i = i\end{math}
\end{itemize}
\end{frame}

\begin{frame}
\frametitle{Alias method: визуальная интерпретация}
\usemintedstyle{solarized-light}
\begin{itemize}
\item У этого метода есть забавная визуальная интерпретация в терминах разрезания и запаковки прямоугольников
\end{itemize}
\slideimage{alias.png}
\end{frame}

\begin{frame}
\frametitle{Alias method}
\usemintedstyle{solarized-light}
\begin{itemize}
\item Alias-таблица для конкретного распределения \textit{не уникальна}
\pause
\item Семплинг работает чуть быстрее, если пытаться максимизировать значения \begin{math}q_i\end{math} (будет меньше обращений к массиву \begin{math}k_i\end{math})
\pause
\item Построение оптимальных массивов \begin{math}q_i\end{math} и \begin{math}k_i\end{math} -- \textbf{NP-трудная} задача
\pause
\item Можно использовать жадный алгоритм, всегда выбирая \begin{math}i\in\operatorname{small}\end{math} с наибольшим \begin{math}q_i\end{math} и \begin{math}j\in\operatorname{large}\end{math} с наименьшим \begin{math}q_j\end{math} -- для этого нужны приоритетные очереди, и время предобработки вырастает до \begin{math}O(N \cdot \log N)\end{math}
\pause
\item Очень подробная статья с выводом алгоритма: \href{https://www.keithschwarz.com/darts-dice-coins}{\texttt{keithschwarz.com/darts-dice-coins}}
\end{itemize}
\end{frame}

\begin{frame}
\frametitle{Unidirectional path tracing}
\usemintedstyle{solarized-light}
\begin{itemize}
\item Алгоритм, который мы используем, называется \textit{unidirectional path tracing}: мы генерируем случайный \underline{путь} света в \textit{одном направлении}, т.е. из камеры в сцену
\pause
\item Многие индустриальные рейтрейсеры используют этот алгоритм (напр. движок Cycles в Blender'е)
\end{itemize}
\end{frame}

\begin{frame}
\frametitle{Unidirectional path tracing}
\usemintedstyle{solarized-light}
\begin{itemize}
\item Recap: генерировать путь из источника света -- плохая идея: он вряд ли (с вероятностью 0) попадёт в камеру
\pause
\item Recap: генерировать путь из камеры -- плохая идея: он может не попасть в источник света
\pause
\item Мы решили вторую проблему importance sampling'ом, но он не всегда поможет
\pause
\item Например, источник света закрыт большим числом другой геометрии, или спрятан за сильно неоднородным стеклом (\textit{frosted glass})
\pause
\item \begin{math}\Longrightarrow\end{math} \textit{Bidirectional path tracing (BDPT)}
\end{itemize}
\end{frame}

\begin{frame}
\frametitle{Bidirectional path tracing (Pixar RenderMan)}
\slideimage{bidirectional_1.jpg}
\end{frame}

\begin{frame}
\frametitle{Bidirectional path tracing (Veach, 1997)}
\usemintedstyle{solarized-light}
\begin{itemize}
\item Возьмём area formulation уравнения рендеринга и переформулируем его как интеграл по \textit{всем возможным путям}, которыми свет может пройти от источников света до камеры
\pause
\item Нужно аккуратно формализовать пространство путей: возьмём объединение пространств всех путей конечной длины
\pause
\item Нужно аккуратно формализовать меру на путях фиксированной длины \begin{math}N\end{math}: она берётся как \begin{math}N\end{math}-кратное произведение мер площади на поверхностях объектов
\pause
\item Нужно аккуратно написать сам интеграл: в него будут входить значения BRDF промежуточных точек, а также \textit{geometry terms}
\pause
\item Нужно придумать, как семплить произвольные пути, и как считать их вероятности для Монте-Карло интегрирования
\end{itemize}
\end{frame}

\begin{frame}
\frametitle{Path integral formulation}
\usemintedstyle{solarized-light}
\begin{itemize}
\item Пусть мы сгенерировали путь \begin{math}p_0 p_1 p_2 \dots p_N\end{math}
\pause
\item Из точки \begin{math}p_0\end{math} в направлении точки \begin{math}p_1\end{math} излучилось \begin{math}L_e(p_0 \rightarrow p_1)\end{math} света
\pause
\item Путь \begin{math}p_0 \rightarrow p_1\end{math} мог быть чем-то закрыт, за что отвечает \textit{visibility term} \begin{math}V(p_0\rightarrow p_1)\end{math}
\pause
\item Нужно учесть преобразование координат интегрирования из направлений на сфере в точки на поверхности: \begin{equation*}\frac{\|\cos \theta_o(p_0) \cos\theta_i(p_1)\|}{\|p_1 - p_0\|^2}\end{equation*}
\pause
\item Вместе с \begin{math}V(p_0\rightarrow p_1)\end{math} это даст \textit{geometry term} \begin{math}G(p_0\rightarrow p_1)\end{math}
\pause
\item Нужно учесть BRDF \begin{math}f(p_0\rightarrow p_1 \rightarrow p_2)\end{math} в точке \begin{math}p_1\end{math} для света, пришедшего из \begin{math}p_0\end{math} и уходящего в \begin{math}p_2\end{math}
\end{itemize}
\end{frame}

\begin{frame}
\frametitle{Path integral formulation}
\usemintedstyle{solarized-light}
\begin{itemize}
\item В итоге, количество света, полученного из этого пути, равно
\end{itemize}
\begin{gather*}
L_e(p_0\rightarrow p_1)\cdot G(p_0\rightarrow p_1) \cdot \\
\cdot f(p_0\rightarrow p_1\rightarrow p_2) \cdot G(p_1\rightarrow p_2) \cdot \\
\cdot f(p_1\rightarrow p_2\rightarrow p_3) \cdot G(p_2\rightarrow p_3) \cdot \\
\dots \\
\cdot f(p_{N-2}\rightarrow p_{N-1}\rightarrow p_N) \cdot G(p_{N-1}\rightarrow p_N)
\end{gather*}
\end{frame}

\begin{frame}
\frametitle{Bidirectional path tracing}
\usemintedstyle{solarized-light}
\begin{itemize}
\item Как генерировать случайные пути?
\pause
\item Один из вариантов:
\pause
\begin{itemize}
\item Сгенерировать случайный путь \begin{math}p_0\dots p_S\end{math} из источника света
\pause
\item Сгенерировать случайный путь \begin{math}q_0\dots q_T\end{math} из камеры
\pause
\item Склеить их в путь \begin{math}p_0\dots p_S q_T \dots q_0\end{math} 
\end{itemize}
\pause
\item На каждом шаге генерации пути останавливаемся с некоторой вероятностью \begin{math}p\end{math} (тогда \begin{math}\frac{1}{p}\end{math} -- средняя длина пути)
\pause
\item Генерировать путь можно, пуская лучи из очередной точки в случайном направлении (можно с importance sampling'ом BRDF, например cosine-weighted для диффузных материалов)
\end{itemize}
\end{frame}

\begin{frame}
\frametitle{Bidirectional path tracing}
\usemintedstyle{solarized-light}
\begin{itemize}
\item Вместо того, чтобы обработать путь \begin{math}p_0\dots p_S q_T \dots q_0\end{math} и выбросить, мы можем дополнительно попарно склеить все префиксы пути \begin{math}p_0\dots p_S \end{math} со всеми префиксами \begin{math}q_0\dots q_T\end{math} и получить квадратичное число путей
\pause
\item BDPT позволяет улучшить сходимость Монте-Карло интегрирования, и позволяет обрабатывать гораздо более сложные сцены
\pause
\item BDPT заметно сложнее в реализации
\pause
\item \textbf{\alert{N.B.}}: Существуют сцены, для которых задача решения уравнения рендеринга \textbf{алгоритмически неразрешима} (см. Reif et al. 1994)
\end{itemize}
\end{frame}

\begin{frame}
\frametitle{Bidirectional path tracing (Pixar RenderMan)}
\slideimage{bidirectional_2.jpg}
\end{frame}

\begin{frame}
\frametitle{Bidirectional path tracing (Sohil Shah)}
\slideimage{bidirectional_3.jpg}
\end{frame}

\begin{frame}
\frametitle{Bidirectional path tracing (Sohil Shah)}
\slideimage{bidirectional_4.jpg}
\end{frame}

\begin{frame}
\frametitle{Bidirectional path tracing: ссылки}
\begin{itemize}
\item \href{https://graphics.stanford.edu/papers/veach_thesis/thesis.pdf}{\texttt{Eric Veach PhD thesis, 1997}} (chapters 8-10)
\item \href{https://graphics.stanford.edu/courses/cs348b-03/papers/mc-course.pdf}{\texttt{Stanford raytracing course notes}}
\end{itemize}
\end{frame}

\begin{frame}
\frametitle{Metropolis light transport}
\usemintedstyle{solarized-light}
\begin{itemize}
\item Path integral formulation превращает задачу рендеринга в вычисление интеграла многомерной (вообще говоря, бесконечномерной) функции
\pause
\item Для этой общей задачи известны хорошие методы, такие как алгоритм Метрополиса-Гастингса
\pause
\item Комбинация этого алгоритма с path integral formulation даёт подход, известный как \textit{metropolis light transport}
\end{itemize}
\end{frame}

\begin{frame}
\frametitle{Алгоритм Метрополиса-Гастингса}
\usemintedstyle{solarized-light}
\begin{itemize}
\item Пусть, есть распределение вероятности \begin{math}p(x)\end{math} (на каком-то множестве), которое мы знаем с точностью до постоянного множителя (т.е. не знаем нормировку)
\pause
\item Другими словами, мы знаем некоторую функцию \begin{math}f(x) = C\cdot p(x)\end{math}, но не знаем \begin{math}C\end{math}
\pause
\item Мы хотим генерировать случайные семлы \begin{math}X\end{math} с вероятностью \begin{math}p(x)\end{math}
\end{itemize}
\end{frame}

\begin{frame}
\frametitle{Алгоритм Метрополиса-Гастингса}
\usemintedstyle{solarized-light}
\begin{itemize}
\item Выберем какой-нибудь первый семпл \begin{math}X_0\end{math}
\pause
\item Выберем некоторое распределение \begin{math}g(x \rightarrow y)\end{math} (\textit{proposal function}), которое по текущей точке \begin{math}x\end{math} будет генерировать следующую точку \begin{math}y\end{math}
\pause
\item Proposal function должна быть симметричной: \begin{math}g(x \rightarrow y) = g(y \rightarrow x)\end{math}
\pause
\item Далее в цикле:
\pause
\begin{itemize}
\item Генерируем новый семпл \begin{math}Y\end{math} используя proposal распределение \begin{math}g(X_i \rightarrow Y)\end{math}
\pause
\item Вычисляем \textit{acceptance ratio} \begin{math}\alpha = \frac{f(Y)}{f(X_i)}\end{math}
\pause
\item С вероятностью \begin{math}\min(\alpha, 1)\end{math} положим \begin{math}X_{i+1} = Y\end{math}, иначе \begin{math}X_{i+1} = X_i\end{math}
\end{itemize}
\pause
\item Тогда последовательность семплов \begin{math}\{X_i\}\end{math} в пределе стремится к распределению \begin{math}p(x)\end{math}
\end{itemize}
\end{frame}

\begin{frame}
\frametitle{Алгоритм Метрополиса-Гастингса}
\usemintedstyle{solarized-light}
\begin{itemize}
\item Если \begin{math}f(Y) > f(X_i)\end{math}, то \begin{math}\alpha > 1\end{math}, и новый семпл будет всегда принят \begin{math}X_{i+1} = Y\end{math}
\pause
\item Чем меньше отношение \begin{math}\alpha = \frac{f(Y)}{f(X_i)}\end{math}, тем меньше вероятность, что семпл будет принят
\pause
\item Интуитивно, алгоритм идёт в направлении увеличения веса семплов, иногда заходя в регионы с меньшей вероятностью
\pause
\item Если случайная величина задана на \begin{math}\mathbb R^n\end{math}, в качестве \begin{math}g(x \rightarrow y)\end{math} часто берут Гауссово распределение \begin{math}g(x \rightarrow y) \sim \exp(-\|x-y\|^2)\end{math}
\end{itemize}
\end{frame}

\begin{frame}
\frametitle{Алгоритм Метрополиса-Гастингса}
\usemintedstyle{solarized-light}
\begin{itemize}
\item Алгоритм очень чувствителен к выбору первого семпла \begin{math}X_0\end{math} (\textit{start-up bias}), поэтому часто используют \textit{burn-in period}: большое количество первых семплов выкидывают
\pause
\item Последовательные семплы алгоритма скоррелированы, что может негативно влиять на результат вычислений
\end{itemize}
\end{frame}

\begin{frame}
\frametitle{Метрополис-Гастингс + Монте-Карло интегрирование}
\usemintedstyle{solarized-light}
\begin{itemize}
\item Хотим вычислить интеграл \begin{math}I=\int f(x)dx\end{math} от некоторой функции \begin{math}f(x)\end{math}
\pause
\item Выбираем плотность вероятности \begin{math}p(x)\end{math}, похожую на нашу функцию, которую мы умеем нормировать
\pause
\item Используем алгоритм Метрополиса-Гастингса с плотностью \begin{math}p(x)\end{math}
\pause
\item Делаем burn-in, вхолостую генерируя набор семплов
\pause
\item После чего генерируем \begin{math}N\end{math} семплов
\pause
\item Аппроксимируем интеграл как
\begin{equation*}I \approx \frac{1}{N}\sum \frac{f(X_i)}{p(X_i)}\end{equation*}
\end{itemize}
\end{frame}

\begin{frame}
\frametitle{Метрополис-Гастингс + Монте-Карло + importance sampling}
\usemintedstyle{solarized-light}
\begin{itemize}
\item Хотим вычислить интеграл \begin{math}I=\int f(x)g(x)dx\end{math} от произведения некоторых функций
\pause
\item Выбираем плотность вероятности \begin{math}p(x) \sim f(x)\end{math}, пропорциональную \begin{math}f(x)\end{math}
\begin{equation*}p(x) = \frac{f(x)}{\int f(y)dy}\end{equation*}
\pause
\item Используем алгоритм Метрополиса-Гастингса с плотностью \begin{math}f(x)\end{math}
\pause
\item Делаем burn-in, вхолостую генерируя набор семплов
\pause
\item Генерируем \begin{math}N\end{math} семплов
\pause
\item Аппроксимируем интеграл как
\begin{equation*}I \approx \frac{1}{N}\sum \frac{f(X_i)g(X_i)}{p(X_i)} = \left[\frac{1}{N}\sum g(X_i)\right]\cdot \left[\int f(x) dx\right]\end{equation*}
\end{itemize}
\end{frame}

\begin{frame}
\frametitle{Metropolis light transport}
\usemintedstyle{solarized-light}
\begin{itemize}
\item Нужно придумать, какую использовать плотность вероятности (возможно, пропорциональную какому-то множителю), и учесть то, что мы не знаем нормировки (возможно, её легче почитать через BDPT)
\pause
\item Нужно сгенерировать первый путь (семпл), например, так же, как в BDPT
\pause
\item Нужно придумать proposal distribution -- случайные мутации пути
\pause
\item \href{https://graphics.stanford.edu/papers/metro/metro.pdf}{\texttt{Veach, Guibas - Metropolis Light Transport}}
\item \href{https://www.pbr-book.org/3ed-2018/Light_Transport_III_Bidirectional_Methods/Metropolis_Light_Transport}{\texttt{PBR Book: Metropolis Light Transport}}
\end{itemize}
\end{frame}

\begin{frame}
\frametitle{Формат glTF}
\begin{itemize}
\item Один из самых распространённых сегодня форматов описания трёхмерных сцен
\pause
\item Поддерживает иерархии объектов и их преобразований
\pause
\item Поддерживает текстуры и сложные материалы
\pause
\item Поддерживает анимированные объекты
\pause
\item Поддерживает источники света (в виде расширения)
\pause
\item Удобен для загрузки на GPU для рендеринга (например, через OpenGL или Vulkan)
\end{itemize}
\end{frame}

\begin{frame}
\frametitle{Формат glTF}
\begin{itemize}
\item Особенность формата в том, что он хранит метаинформацию (описание дерева объектов, и т.д.) в текстовом виде (в JSON'е), а бинарные данные (координаты отдельных вершин, индексы вершин, keyframe'ы анимаций) -- в сыром бинарном виде
\pause
\item Бинарные данные обычно представлены ссылкой на внешний источник (файл, URL, etc)
\pause
\item Есть вариация формата (GLB), которая склеивает JSON и бинарные данные в один файл
\pause
\item \href{https://registry.khronos.org/glTF/specs/2.0/glTF-2.0.html}{\texttt{Ссылка на спецификацию формата}}
\end{itemize}
\end{frame}

\begin{frame}
\frametitle{Формат glTF: структура}
\usemintedstyle{solarized-light}
\begin{itemize}
\item Объект формата glTF -- это JSON-объект, большинство полей которого это массивы объектов определённого \textit{типа}: nodes, meshes, accessors, buffers, bufferViews, materials, cameras, ...
\pause
\item Объекты ссылаются друг на друга с помощью индексов в эти массивы
\pause
\item Ноды образуют виртуальное дерево объектов сцены (необязательно видимых объектов), и могут содержать другие ноды в качестве детей
\pause
\item Каждая нода может иметь матрицу преобразования относительно родительской ноды (или относительно центра сцены, если это корневая нода)
\end{itemize}
\end{frame}

\begin{frame}
\frametitle{Формат glTF: структура}
\usemintedstyle{solarized-light}
\begin{itemize}
\item Нода может иметь меш (тогда эта нода подлежит рисованию)
\pause
\item Меш содержит несколько примитивов, каждый из которых имеет свой материал и набор вершин
\pause
\item Позиции вершин и индексы вершин в треугольниках описываются \textit{аксессорами}
\pause
\item Аксессоры задают тип и формат данных, и ссылаются на buffer view, которые ссылаются на буфер с бинарными данными
\end{itemize}
\end{frame}

\begin{frame}
\frametitle{Формат glTF: структура}
\usemintedstyle{solarized-light}
\begin{itemize}
\item Нода может содержать камеру -- тогда позиция и ориентация этой ноды задают позицию и ориентацию камеры
\pause
\item Вид и параметры проекции камеры задаются соответствующим объектом из массива cameras
\end{itemize}
\end{frame}

\begin{frame}
\frametitle{Формат glTF: реализация}
\usemintedstyle{solarized-light}
\begin{itemize}
\item Для парсинга формата нужен какой-нибудь парсер JSON
\pause
\item C++: Rapidjson, nlohmann/json
\pause
\item Rust: json-rust
\end{itemize}
\end{frame}

\begin{frame}
\frametitle{Формат glTF: реализация}
\usemintedstyle{solarized-light}
\begin{itemize}
\item Для использования в нашем рейтрейсере, нам нужно превратить glTF-сцену в
\pause
\begin{itemize}
\item Набор вершин
\pause
\item Набор треугольников
\pause
\item Набор материалов
\pause
\item Камеру
\end{itemize}
\end{itemize}
\end{frame}

\begin{frame}
\frametitle{Формат glTF: реализация}
\usemintedstyle{solarized-light}
\begin{itemize}
\item В качестве камеры берём первую попавшуюся ноду с камерой (считаем, что такая нода ровно одна)
\pause
\item Поддерживаем только перспективную камеру
\pause
\item Позицию и ориентацию камеры берём из содержащей её ноды
\end{itemize}
\end{frame}

\begin{frame}
\frametitle{Формат glTF: реализация}
\usemintedstyle{solarized-light}
\begin{itemize}
\item Треугольник описывается тройкой индексов вершин и индексом материала
\pause
\item Можно хранить разные объекты (ноды/меши/примитивы) отдельно, а можно слить всё в одну кучу треугольников
\pause
\item В любом случае, нужно учесть преобразование ноды -- его можно применить к вершинам сразу, а можно применять обратное преобразование непосредственно при поиске пересечения
\pause
\item Можно построить один BVH на все треугольники, а можно разбить на отдельные TLAS + BLAS
\end{itemize}
\end{frame}

\begin{frame}
\frametitle{Формат glTF: реализация}
\usemintedstyle{solarized-light}
\begin{itemize}
\item Преобразования нод могут быть заданы или явной матрицей \begin{math}4\times 4\end{math} (в однородных координатах), или тройкой translation-rotation-scale
\pause
\item Translation задаётся трёхмерным вектором
\pause
\item Rotation задаётся единичным кватернионом
\pause
\item Scale задаётся трёхмерным вектором и может описывать неизотромное масштабирование (\textit{nonuniform scaling}) -- его нужно покомпонентно умножать на координаты вершины
\pause
\item К преобразованию ноды нужно рекурсивно добавить преобразования всех родительских нод
\end{itemize}
\end{frame}

\begin{frame}
\frametitle{Формат glTF: реализация}
\usemintedstyle{solarized-light}
\begin{itemize}
\item От самих вершин нам пока нужны только координаты (не забываем к ним применить пребразование ноды)
\pause
\item Они могут быть заданы в большом количестве разных форматов; будем считать, что они всегда в виде трёхмерных 32-битных floating point
\pause
\item Также будем считать, что индексы или 16-битные, или 32-битные беззнаковые целые числа (придётся поддержать два варианта)
\pause
\item Обратите внимание, что byteOffset есть и в аксессорах, и в bufferView
\pause
\item Параметр byteStride в bufferView будем игнорировать
\end{itemize}
\end{frame}

\begin{frame}
\frametitle{Формат glTF: реализация}
\usemintedstyle{solarized-light}
\begin{itemize}
\item glTF поддерживает довольно сложные материалы; пока мы будем использовать только несколько параметров
\pause
\item \mintinline{cpp}|pbrMetallicRoughness.baseColorFactor| -- цвет материала (RGBA)
\pause
\item Если альфа-канал цвета меньше 1, считаем материал диэлектриком с IOR = 1.5
\pause
\item Если \mintinline{cpp}|pbrMetallicRoughness.metallicFactor| больше 0, считаем объект металлом
\pause
\item Иначе объект диффузный
\pause
\item В любом случае, \mintinline{cpp}|emissiveFactor| -- интенсивность излучения материала
\end{itemize}
\end{frame}

\begin{frame}
\frametitle{Формат glTF: реализация}
\usemintedstyle{solarized-light}
\begin{itemize}
\item Обратите внимание на многочисленные дефолты: например, \mintinline{cpp}|baseColorFactor| по умолчанию равен \mintinline{cpp}|[1, 1, 1, 1]|, а \mintinline{cpp}|metallicFactor| равен 1
\pause
\item По спецификации, \mintinline{cpp}|emissiveFactor| не может быть больше 1; это обходится расширением \mintinline{cpp}|KHR_materials_emissive_strength| (см. примеры сцен для практики)
\end{itemize}
\end{frame}

\end{document}
